\begin{abstract}
Graphs have become a common language for artificial intelligence, from social
and biological interaction networks to knowledge graphs, retrieval pipelines,
and autonomous-agent workflows.
A central challenge in these systems is to identify meaningful organization
across scales rather than only local interactions.
\emph{Structural Entropy} (SE) addresses this challenge by extending
information-theoretic reasoning from pairwise uncertainty to hierarchical
structure: it measures how efficiently a graph can be encoded by a multilevel
partition and, through minimization, reveals candidate structural abstractions
of the system.

Over the past decade, SE has developed from a community-detection principle
into a broader framework spanning graph clustering, graph pooling, graph
structure learning, reinforcement learning, bioinformatics, social-network
analysis, and pattern recognition~\cite{Li2016,wu2022pooling,zeng2025sidmJMLR,zhang2021supertad,peng2024zjyun,zhu2023HiTIN}.
In this survey we review that development.
We summarize the mathematical foundations of SE, organize the literature into
major methodological families, compare SE with neighboring graph objectives
such as modularity and flow-based compression, and examine how the field is
evaluated in practice.
We also discuss open problems in scalability, differentiable optimization,
heterogeneous structures, and integration with modern graph-enhanced AI
systems.
A companion open-source library, \texttt{glass-jax}, accompanies this survey
as a reproducibility artifact.

\medskip\noindent\textbf{Keywords:} structural entropy, graph clustering,
hierarchical community detection, graph neural networks, reinforcement learning,
information theory.
\end{abstract}
